\documentclass[10pt,a4paper]{amsart}
\usepackage[margin=19mm]{geometry}
\usepackage{amsmath,amssymb,mathtools}
\usepackage[scr=rsfs,cal=euler]{mathalfa}
\usepackage{xfrac}
\usepackage[unicode,hidelinks,bookmarks=false]{hyperref}
\newcommand{\ie}{i.e.\ }
\newcommand{\cf}{cf.\ }
\newcommand{\resp}{respectively\ }
\newcommand{\Subsection}{\subsection}
\providecommand{\nobreakdash}{\nobreak\hbox{-}}
\setlength{\emergencystretch}{2em}
\multlinegap0pt
\usepackage{amssymb}
\usepackage{mathtools}
\usepackage{bm}
\newtheorem{corollary}{Corollary}
\newtheorem{assumption}{Assumption}
\newtheorem{theorem}{Theorem}
\newcommand{\ms}[1]{(\mathbf{#1},\mathcal{#1})}
\title{The Variational Approach in Filtering and Correlated Noise}
\author{Sharan Srinivasan, Vijay Gupta, Harsha Honnappa}
\date{Source: arXiv:2604.03001v1 (CC BY 4.0)}
\begin{document}
\maketitle

\noindent\emph{Source and licence.} The Variational Approach in Filtering and Correlated Noise. Version arXiv:2604.03001v1 (CC BY 4.0), distributed by arXiv under the Creative Commons Attribution 4.0 International licence, http://creativecommons.org/licenses/by/4.0/, as recorded in the arXiv licence metadata for this exact version and shown on the abstract page https://arxiv.org/abs/2604.03001v1. Full licence notice: \texttt{licenses/arxiv-2604.03001-CC-BY-4.0.txt}.\par\smallskip

\noindent\textit{Editorial scope.} Counted unit: the article's theorem theorem (source0.tex lines 332--334). The article's complete original proof of it is delivered with it (source0.tex lines 335--416, one \texttt{proof} environment of 82 lines). No mathematics is written, repaired or re-derived: the statement and the proof are the article's own lines, in the article's own notation, and no retained line is substituted or re-flowed.  Retained uncounted as the source-local context the proof needs: lines 221--227; lines 246--251; lines 322--328.  Nothing was omitted from any retained span, so the package carries no omission record.  No retained line contains a citation, so the wrapper carries no bibliography and no bibliography entry.  No retained line contains a figure environment, an image-inclusion command or drawing code, so no image is delivered and the package declares no asset dependency.  Editorial formatting only, in addition: an AMS \texttt{amsart} preamble that restates the article's own declarations and macros used by the retained lines, namely the environments \texttt{corollary}, \texttt{assumption}, \texttt{theorem} on separate counters, together with the standard packages those lines need.  The retained environments print in this document as Corollary 1, Assumption 1, Theorem 1 in order of appearance, whereas the article prints the same items with the numbering of its own full document.  The complete native source of the article as distributed in the arXiv e-print for this exact version is bundled unchanged as \texttt{sources/197778/source0.tex}.
\begin{corollary}\label{cor:abstract}
    Let $P_{X|Y}(\cdot, y)$ denote a regular conditional distribution of $X$ given $Y$. If
    \[
        P_Y\!\left(\left\{ y \in \mathcal{Y} : P_{X|Y}(\cdot, y) \perp P_X \right\}\right) > 0,
    \]
    then there is no $\sigma$-finite measure $\lambda_Y$ on $\ms{Y}$ satisfying Assumption~2.1.
\end{corollary}

\begin{align}
    \label{eq:sig corr}
    X_t &= X_0 + \int_0^tb(s,X_s)ds+\int_0^t\sigma_0(s,X_s)dB_s+\int_0^t\sigma_1(s,X_s)dW_s,\\
    \label{eq:obs corr}
    Y_t &= \int_0^th(X_s)ds + W_t,
\end{align}

\begin{assumption}\label{as:hordmander's condition}
    The set
    \[
    Z := \{z \in \mathbb R^d: \sigma_1(t,z)=0\},
    \]
    has Lebesgue measure 0 a.e.\ $t\in[0,T]$, and the vector fields $\sigma_0$ satisfy the H\"ormander condition at $X_0$ almost surely.
\end{assumption}

\begin{theorem}
    Given SDEs (\ref{eq:sig corr}), (\ref{eq:obs corr}), with joint measure $P_{XY}$ and Assumption \ref{as:hordmander's condition} is satisfied. Then there is no measure $\lambda_Y$ on $\ms{Y}$ such that $P_{XY}\ll P_X\otimes \lambda_Y$.
\end{theorem}

\begin{proof}
    By Corollary~\ref{cor:abstract}, it suffices to show that $P_{XY} \perp P_X\otimes P_Y$. Consider the map $Q_n^{[0,t]}:\mathbf{X}\times\mathbf{Y}\to \mathbb{R}^{d\times n}$ defined as:
    \begin{equation}
        Q_n^{[0,t]}(x,y) = \sum_{i=0}^{2^n-1}(x_{t_{i+1}}-x_{t_i})(y_{t_{i+1}}-y_{t_i})^T,
    \end{equation}
    where the sum is taken over dyadic intervals of $[0,t]\subseteq[0,T]$.
    Under the joint measure $P_{XY}$, we have the quadratic variation:
    \begin{align*}
        \lim_{n\to\infty}Q_n^{[0,t]}(X,Y)  &= \lim_{n\to\infty}\sum_{i=0}^{2^n-1}(X_{t_{i+1}}-X_{t_i})(Y_{t_{i+1}}-Y_{t_i})^T\\
        &= \langle X,Y\rangle_t = \int_0^t\sigma_1(s,X_s)ds,
    \end{align*}
    almost surely.

    Under the product measure $P_X\otimes P_Y$, the mean of $Q_n^{[0,t]}$ under the limit $n\to\infty$ is
    \begin{align*}
        \mathbb{E}[Q_n^{[0,t]}(X,Y)] &= \sum_{i=0}^{2^n-1} \mathbb{E}_X[X_{t_{i+1}}-X_{t_i}]\mathbb{E}_Y[Y_{t_{i+1}}-Y_{t_i}]^T \\
         &= \sum_{i=0}^{2^n-1}\int_{t_i}^{t_{i+1}}\mathbb{E}_X\left[b(s,X_s)\right]ds\times\int_{t_i}^{t_{i+1}}\mathbb{E}_Y\left[h(X_s)\right]^Tds.
    \end{align*}
    Since both $b$ and $h$ have at most linear growth, we have:
    \begin{align*}
        \mathbb{E}[|Q_n^{[0,t]}(X,Y)|] &\leq C\sum_{i=0}^{2^n-1}\int_{t_i}^{t_{i+1}}\mathbb{E}_X[1+|X_s|]ds\times\int_{t_i}^{t_{i+1}}\mathbb{E}_Y[1+|Y_s|]ds\\
        &\leq \tilde{C}\sum_{i=0}^{2^n-1}(t_{i+1}-t_i)^2\to0,
    \end{align*}
    where the constant $\tilde{C} := C\sup_{0\leq s\leq T}\mathbb{E}_X[1+|X_s|]\mathbb{E}_Y[1+|Y_s|]\}$ is finite, since the moments of a strong solution of an It\^o SDE are bounded. Hence $\lim_{n\to\infty}\mathbb{E}[Q_n^{[0,t]}(X,Y)] = 0$ under the product measure.

    The variance under the limit $n\to\infty$ is
    \begin{align*}
        \text{Var}(Q_n^{[0,t]}) &= \sum_{i=0}^{2^n-1} \mathbb{E}_X[(X_{t_{i+1}}-X_{t_i})^2]\mathbb{E}_Y[(Y_{t_{i+1}}-Y_{t_i})^2] \\
        &= \sum_{i=0}^{2^n-1}\int_{t_i}^{t_{i+1}}\mathbb{E}_X\left[(|\sigma_0(s,X_s)|^2+|\sigma_1(s,X_s)|^2)\right]ds\times\int_{t_i}^{t_{i+1}}ds.
    \end{align*}
    Since the maps $\sigma_0$ and $\sigma_1$ have at most linear growth,
    \begin{align*}
        \text{Var}(Q_n^{[0,t]}) &\leq C'\sum_{i=0}^{2^n-1}\int_{t_i}^{t_{i+1}}\mathbb{E}_X[(1+|X_s|)^2]ds\times (t_{i+1}-t_i)\\
        &\leq \tilde{C}'\sum_{i=0}^{2^n-1}(t_{i+1}-t_i)^2\to 0,
    \end{align*}
    where $\tilde{C}' := C'\sup_{0\leq s\leq T}\mathbb{E}_X[(1+|X_s|)^2]$ is bounded due to bounded moments of a strong solution to an It\^o SDE.

    Thus, we have $Q_n^{[0,t]}\to 0$ in $L^2$ as $n\to\infty$. From Markov's inequality, it follows that $Q_n^{[0,t]}(X,Y)\to0$ in probability as $n\to\infty$. In order to show convergence almost surely, we choose an increasing subsequence $\{n_k\}_{k\in \mathbb{N}}$ such that
    \begin{equation*}
        \mathbb{P}(|Q^{[0,t]}_{n_k}(X,Y)|>1/k) \leq \frac{1}{2^k}.
    \end{equation*}
    Since
    \[
        \sum_{k\in\mathbb{N}}\mathbb{P}(|Q^{[0,t]}_{n_k}(X,Y)|>1/k)  \leq 1,
    \]
    from the Borel-Cantelli lemma,
    \begin{equation*}
        \mathbb{P}(\limsup_{k\to\infty}\{|Q^{[0,t]}_{n_k}(X,Y)|>1/k\}) = 0,
    \end{equation*}
    i.e., the event $\{|Q^{[0,t]}_{n_k}(X,Y)|>1/k\}$ occurs only finitely many times almost surely.
    Therefore, for almost every $\omega\in\Omega$, there exists $K(\omega)$ such that for all $l>K(\omega)$,
    \[
        |Q^{[0,t]}_{n_l}(X,Y)|\leq\frac{1}{l}.
    \]
    Since the sequence $\{n_k\}$ is increasing, we have $Q_n^{[0,t]}(X,Y) \to 0$ almost surely as $n\to\infty$.
    Define the sets
    \begin{align*}
        A^{[0,t]} &= \left\{(x,y)\in\mathbf{X\times Y}:\lim_{n\to\infty}Q_n^{[0,t]}(x,y) = \int_0^t\sigma_1(s,x_s)ds\right\},\\
        B^{[0,t]} &= \{(x,y)\in\mathbf{X\times Y}:\lim_{n\to\infty}Q_n^{[0,t]}(x,y) =0\}.
    \end{align*}
    We have $P_{XY}(A^{[0,t]}) = 1$ and $(P_X\otimes P_Y)(B^{[0,t]}) =1$ for all $t\in [0,T]$. In order to show the mutual singularity of the two measures, we need to show that for some $s\in[0,T]$, $(P_X\otimes P_Y)(A^{[0,s]}) = 0$. For all $t\in[0,T]$, we have
    \begin{align*}
        (P_X\otimes P_Y)(A^{[0,t]}) &= (P_X\otimes P_Y)(A^{[0,t]}\cap B^{[0,t]}) + (P_X\otimes P_Y)(A^{[0,t]}\cap (B^{[0,t]})^c).
    \end{align*}
    Since $(P_X\otimes P_Y)((B^{[0,t]})^c) = 0$, $(P_X\otimes P_Y)(A^{[0,t]}\cap (B^{[0,t]})^c) = 0$. The set $A^{[0,t]}\cap B^{[0,t]}$ is given by:
    \begin{equation*}
        A^{[0,t]}\cap B^{[0,t]} = \left\{(x,y)\in\mathbf{X\times Y}:\int_0^t\sigma_1(s,x_s)ds = 0\right\}.
    \end{equation*}
    If there exists a $t\in[0,T]$ such that $A^{[0,t]}\cap B^{[0,t]} = \emptyset$, then we have $(P_X\otimes P_Y)(A^{[0,t]}) = (P_X\otimes P_Y)(A^{[0,t]}\cap B^{[0,t]}) = 0$, we are done. If there is no $t\in[0,T]$ such that $A^{[0,t]}\cap B^{[0,t]} = \emptyset$, then there exists $x\in\mathbf{X}$ such that
    \[
        \int_0^t\sigma_1(s,x_s)ds = 0,
    \]
    for all $t\in[0,T]$. This means that $\sigma_1(t,x_t) = 0$ for almost every $t\in[0,T]$. We now show that the product measure $P_X\otimes P_Y$ assigns no mass to the set
    \[
        S := \{(x,y)\in\mathbf{X\times Y}: \sigma_1(t,x_t) = 0 \text{ a.e.\ } t\in[0,T]\}.
    \]
    It is enough to show that $P_X(\{x:\sigma_1(t,x_t) = 0 \text{ a.e.\ } t\in[0,T]\}) = 0$. From Assumption \ref{as:hordmander's condition}, $X_t$ admits a density with respect to the Lebesgue measure. Since, $Z = \{z\in\mathbb{R}^d:\sigma_1(t,z) = 0\}$ has Lebesgue measure 0 a.e.\ $t\in[0,T]$, we have
    \begin{align*}
\mathbb{E}_X\left[\int_0^T\mathbf{1}_{\{\sigma_1(t,X_t)=0\}}dt\right] &= \int_0^TP_X(\{x:\sigma_1(t,x_t)=0\})dt\\ &=0.
    \end{align*}
    Thus, $P_X(\{x:\sigma_1(t,x_t)=0\}) = 0$ a.e.\ $t\in[0,T]$, from which we have that $P_X \otimes P_Y$ assigns no mass to the set $S$. Therefore, $P_{XY}$ is not absolutely continuous with respect to $P_X \otimes P_Y$.
\end{proof}

\end{document}
